\documentclass[12pt]{article}
\usepackage{graphicx}
\usepackage[margin=1in, a4paper]{geometry}
\usepackage{setspace}
\usepackage{pdfpages}
\usepackage{float}
\usepackage{amsmath}
\usepackage{fancyvrb}
\usepackage{amssymb}
\usepackage{minted}
\usepackage{enumitem}
\usepackage{tikz-cd}
\usepackage[colorlinks,linkcolor=blue]{hyperref}


\renewcommand{\contentsname}{Contents}
\renewcommand{\figurename}{Figure}
\renewcommand{\tablename}{Table}
\hypersetup{
    colorlinks=true,
    linkcolor=black,
    filecolor=magenta,      
    urlcolor=blue,
}
\newcommand{\mytitle}{Abstract Algebra II Homework 6}
\newcommand{\myauthor}{B13902022 賴昱錡}

\usepackage{fancyhdr}
\pagestyle{fancy}
\fancyhead{}
\fancyhead[L]{\mytitle}
\fancyhead[R]{\myauthor}

\usepackage[english]{babel}
\newtheorem{theorem}{Theorem}[section]
\newtheorem{corollary}{Corollary}[theorem]
\newtheorem{lemma}[theorem]{Lemma}

\usepackage{amsthm}
\theoremstyle{definition}
\newtheorem{definition}{Definition}[section]
\theoremstyle{remark}
\newtheorem*{remark}{Remark}

\usepackage{xeCJK}
\setCJKmainfont{Noto Serif CJK TC}

\title{\mytitle}
\author{\myauthor}
\date{\today}

\begin{document}
\maketitle

\section{} % Problem 1
Let $L$ be the splitting field of $f$ over $K$, and let $G = \mathrm{Gal}(L/K)$. By its action on the four distinct roots $u_1, u_2, u_3, u_4$, we can view $G$ as a subgroup of the symmetric group $S_4$.

The given elements $\alpha, \beta$, and $\gamma$ are the roots of the resolvent cubic of $f$. Since $K$ is a field of characteristic zero, the field extension $K(\alpha, \beta, \gamma)/K$ is Galois, which means:
\[
\mathrm{Aut}(K(\alpha, \beta, \gamma)/K) = \mathrm{Gal}(K(\alpha, \beta, \gamma)/K)
\]

The group $G \le S_4$ acts on the set $X = \{\alpha, \beta, \gamma\}$ by permuting the indices of the roots $u_i$. This action defines a natural group homomorphism from $G$ to the symmetric group on three letters:
\[
\phi: G \to S_3
\]

The image of $\phi$ is exactly the Galois group of $K(\alpha, \beta, \gamma)$ over $K$. To find the structure of this group, we must determine the kernel of $\phi$, which consists of the permutations in $G$ that fix $\alpha, \beta$, and $\gamma$ simultaneously. 

Evaluating the stabilizer of each element in $S_4$:
\begin{itemize}
    \item The stabilizer of $\alpha$ is a dihedral group of order 8: $D_8^{(1)} = \langle(12), (1324)\rangle$.
    \item The stabilizer of $\beta$ is $D_8^{(2)} = \langle(13), (1234)\rangle$.
    \item The stabilizer of $\gamma$ is $D_8^{(3)} = \langle(14), (1243)\rangle$.
\end{itemize}

The intersection of these three stabilizers is the normal Klein four-group in $S_4$:
\[
V_4 = \{ \mathrm{id}, (12)(34), (13)(24), (14)(23) \}
\]

Therefore, the kernel of the homomorphism $\phi$ is $G \cap V_4$. By the First Isomorphism Theorem, the Galois group of the resolvent cubic is isomorphic to the quotient:
\[
\mathrm{Aut}(K(\alpha, \beta, \gamma)/K) \cong G / (G \cap V_4)
\]

Because $V_4 \triangleleft S_4$ and $S_4/V_4 \cong S_3$, the structure of $\mathrm{Aut}(K(\alpha, \beta, \gamma)/K)$ is isomorphic to a subgroup of $S_3$. The exact structure depends on the specific Galois group $G$ of the original quartic $f$:
\begin{itemize}
    \item If $G = S_4$, then $\mathrm{Aut}(K(\alpha, \beta, \gamma)/K) \cong S_3$.
    \item If $G = A_4$, then $\mathrm{Aut}(K(\alpha, \beta, \gamma)/K) \cong \mathbb{Z}/3\mathbb{Z}$.
    \item If $G = D_8$ or $G = \mathbb{Z}/4\mathbb{Z}$, then $\mathrm{Aut}(K(\alpha, \beta, \gamma)/K) \cong \mathbb{Z}/2\mathbb{Z}$.
    \item If $G = V_4$, then $\mathrm{Aut}(K(\alpha, \beta, \gamma)/K)$ is the trivial group.
\end{itemize}

\section{} % Problem 2
\subsection*{(a)}
By the rational root theorem, the possible rational roots are $\pm 1$. But $P(1)=P(-1)=-28\ne 0$, so we know $P(t)$ has no linear factors over $\mathbb{Q}[t]$.

Assume that $P(t)$ factors into two monic quadratics:
\begin{align*}
    P(t)&=(x^2+ax+b)(x^2+cx+d) \\
    &=x^4+(a+c)x^3+(b+d+ac)x^2+(ad+bc)x+bd \\
\end{align*}

By comparing the coefficients, we know:
\begin{align*}
    a+c=0 \\
    b+d+ac=-30 \\
    ad+bc=0 \\
    bd=1
\end{align*}

If $a=c=0$, then $b+d=-30,bd=1$. $b,d$ are the roots for $y^2+30t+1$, clearly, $b,d\notin \mathbb{Q}$. If $c=-a\ne 0$, then $b=d$, and $bd=1$ implies $b=d=\pm 1$. In either cases, $a$ can't be a rational number. 

Therefore no factorization for $P(t)$ exists in $\mathbb{Q}[t]$, so $P(t)$ is irreducible over $\mathbb{Q}$.

\subsection*{(b)}
The all roots of $P(t)=0$ are $\pm 2\sqrt{2}\pm \sqrt{7}$. By grouping the roots, there are 3 possible factorization (with two monic quadratics factors):
\begin{align*}
    P(t)&=(t^2+4\sqrt{2}t+1)(t^2-4\sqrt{2}t+1) \\
    &=(t^2+2\sqrt{7}t-1)(t^2-2\sqrt{7}t-1) \\
    &=(t^2-15+4\sqrt{14})(t^2-15-4\sqrt{14})
\end{align*}

We can find that for every prime $p$ (except 2 and 7), at least one of 2, 7, 14 is a quadratic residue, because when we use Legendre symbols:
$$\left(\frac{2}{p}\right)\left(\frac{7}{p}\right)=\left(\frac{14}{p}\right)$$
If both $\left(\frac{2}{p}\right)=\left(\frac{7}{p}\right)=-1$, then $\left(\frac{14}{p}\right)=1$.

Therefore, if $\exists s\in \mathbb{F}_p$ ($p=\ne2,7$) such that $s^2=2$, then we can write $P(t)=(t^2+4st+1)(t^2-4st+1)$. If $\exists u\in \mathbb{F}_p$ ($p=\ne2,7$) such that $u^2=7$, then we can write $P(t)=(t^2+2ut-1)(t^2-2ut-1)$. If $\exists v\in \mathbb{F}_p$ ($p=\ne2,7$) such that $v^2=14$, then we can write $P(t)=(t^2-15+4v)(t^2-15-4v)$. Here, $P(t)$ is considered under $\mathbb{F}_p[t]$ for some prime $p\ne 2,7$.

Let's consider the cases when $p=2,7$. When $p=2$, $P(t)=x^4+1=(x+1)^{4}\in \mathbb{F}_2[t]$. When $p=7$, $P(t)=x^4+5x^2+1=(x^2-1)^2\in \mathbb{F}_7[t]$.

Hence, $P(t)\mod p\in \mathbb{F}_p[t]$ is reducible for all prime numbers $p$.

\subsection*{(c)}
The four roots of $f(x)$ are $\pm 2\sqrt{2} \pm \sqrt{7}$. Since all roots can be generated by adjoining $\sqrt{2}$ and $\sqrt{7}$ to $\mathbb{Q}$, the splitting field is:
$$K = \mathbb{Q}(\sqrt{2}, \sqrt{7})$$

Since $\sqrt{2}$ and $\sqrt{7}$ are square roots of distinct prime numbers, the degree of the extension is $[K:\mathbb{Q}] = 4$. The Galois group $G = \text{Gal}(K/\mathbb{Q})$ has order 4. Because $K$ is generated by two independent square roots, $G$ is isomorphic to the Klein four-group, $V_4 \cong \mathbb{Z}_2 \times \mathbb{Z}_2$.

The group $G$ consists of four automorphisms, completely determined by their action on the generators $\sqrt{2}$ and $\sqrt{7}$:
\begin{itemize}
    \item $id$: $\sqrt{2} \mapsto \sqrt{2}, \quad \sqrt{7} \mapsto \sqrt{7}$
    \item $\sigma$: $\sqrt{2} \mapsto -\sqrt{2}, \quad \sqrt{7} \mapsto \sqrt{7}$
    \item $\tau$: $\sqrt{2} \mapsto \sqrt{2}, \quad \sqrt{7} \mapsto -\sqrt{7}$
    \item $\sigma\tau$: $\sqrt{2} \mapsto -\sqrt{2}, \quad \sqrt{7} \mapsto -\sqrt{7}$
\end{itemize}

By the Fundamental Theorem of Galois Theory, the subgroups of $G$ correspond to the intermediate fields of $K/\mathbb{Q}$. The proper, non-trivial subgroups of $G$ are $\langle \sigma \rangle$, $\langle \tau \rangle$, and $\langle \sigma\tau \rangle$, all of which have order 2 (and thus index 2 in $G$). We find their corresponding fixed fields:
\begin{itemize}
    \item $K^{\langle \sigma \rangle}$: The automorphism $\sigma$ fixes $\sqrt{7}$. Thus, the intermediate field is $\mathbb{Q}(\sqrt{7})$.
    \item $K^{\langle \tau \rangle}$: The automorphism $\tau$ fixes $\sqrt{2}$. Thus, the intermediate field is $\mathbb{Q}(\sqrt{2})$.
    \item $K^{\langle \sigma\tau \rangle}$: The automorphism $\sigma\tau$ negates both $\sqrt{2}$ and $\sqrt{7}$. Their product, $\sqrt{2}\sqrt{7} = \sqrt{14}$, is mapped to $(-\sqrt{2})(-\sqrt{7}) = \sqrt{14}$. Thus, it fixes $\sqrt{14}$, giving the intermediate field $\mathbb{Q}(\sqrt{14})$.
\end{itemize}
\begin{figure}[H]
    \centering
    % Subgroup Lattice
    \begin{tikzcd}[row sep=4em, column sep=2em]
        & G = \text{Gal}(K/\mathbb{Q}) \arrow[dl, dash, "2"'] \arrow[d, dash, "2"] \arrow[dr, dash, "2"] & \\
        \langle \sigma \rangle \arrow[dr, dash, "2"'] & \langle \sigma\tau \rangle \arrow[d, dash, "2"] & \langle \tau \rangle \arrow[dl, dash, "2"] \\
        & \{id\} &
    \end{tikzcd}
    \hspace{2cm} % Space between the two lattices
    % Intermediate Field Lattice
    \begin{tikzcd}[row sep=4em, column sep=2em]
        & K = \mathbb{Q}(\sqrt{2}, \sqrt{7}) \arrow[dl, dash, "2"'] \arrow[d, dash, "2"] \arrow[dr, dash, "2"] & \\
        \mathbb{Q}(\sqrt{7}) \arrow[dr, dash, "2"'] & \mathbb{Q}(\sqrt{14}) \arrow[d, dash, "2"] & \mathbb{Q}(\sqrt{2}) \arrow[dl, dash, "2"] \\
        & \mathbb{Q} &
    \end{tikzcd}
\end{figure}

\end{document}
% how to display codes?
% \begin{minted}[frame=lines,framesep=2mm,baselinestretch=1.2,linenos,breaklines]{python}
% \end{minted}

% how to display images?
% \begin{figure}[H]
%     \centering
%     \includegraphics[width=0.5\linewidth]{}
%     \caption{Caption}
% \end{figure}
% test
